
Let $(\bm{\mu},\lambda) \in \mathbb{R}_+^{\abs{\cP}} \times \mathbb{R}_+$ be an optimal dual solution to RMP, where dual values $\bm{\mu}$, $\lambda$ correspond to constraints \eqref{eq:BM-misP} and \eqref{eq:BM-complex} respectively.
The reduced cost of a variable $w_k$, associated with a rule $k$ not yet in the master problem, can be expressed as
\begin{equation}
    \hat \rho_k =  \sum_{i \in \cN} \Ind(k \in \K_i)
    -\sum_{i \in \cP}  \mu_i  \Ind(k \in \K_i) +\lambda  c_k.  \label{eq:BP_redcost}
\end{equation}
For the derivation we refer to \autoref{sec:reduced_costs_boolean_model}.
If $\rho_k < 0$, then adding rule $k$ to the relaxed master problem can improve its objective.
If $\rho_k \geq 0$ for all $k \in \K$, we have found an optimal LP solution.
Deciding whether this is the case is called the pricing problem.

To solve the pricing problem,
we construct an IP.
Introduce the variables $z_l\in\{0,1\}, \; l\in \L$ representing whether the rule includes clause $l$.
We also introduce the variables $\delta_i\in\{0,1\}, \; i\in \I$ denoting whether point $i$ is covered by the rule.
The complexity of a rule can be expressed as
\begin{equation}\label{eq:c_z}
    c(\bm{z}) = 1 + \sum_{l \in \L} z_l, \tquad \bm{z} \in \{0,1\}^{\L},
\end{equation}
which we limit by $D$.
Define
\begin{equation}\label{eq:S_i_boolean}
    S_i = \{l \in \L \fw X_{ij} = 0\}, \tquad i \in \I.
\end{equation}
Hence, the set $S_i$ contains all clauses (binary features in this setting) that point $i$ does not satisfy (possess).
The pricing problem may be solved by minimizing \eqref{eq:BP_redcost}.
The corresponding IP, as described in~\cite{lawless_interpretable_2023}, is given by
\begin{subequations}
    \label{eq:BP}
    \begin{align}
        \min \quad & \sum_{i \in \cN} \delta_i - \sum_{i \in \cP} \mu_i \delta_i +
        \lambda c(\bm{z}) \label{eq:BP-obj} \\
        \st \quad
        & D\delta_i + \sum_{j \in S_i} z_j \leq D, && i \in \I, & \label{eq:BP-upper}\\
        & \delta_i +  \sum_{j \in S_i} z_j \geq  1, && i \in \I,  &\label{eq:BP-lower}\\
        & \sum_{j \in \J} z_j \leq D, \quad \label{eq:BP-complex}\\
        & z\in \{0,1\}^{\J},~ \delta\in \{0,1\}^{\I}.\label{eq:BP-int}\qquad
    \end{align}
\end{subequations}

If $\delta_i = 1$, then by \eqref{eq:BP-upper}, no clause that cuts off $i$ can be chosen, so $i$ must covered by the rule.
If $\delta_i = 0$, then by \eqref{eq:BP-lower}, some clause that cuts off $i$ must be chosen, so $i$ cannot be not covered by the rule.
Hence, $\delta_i$ reflects whether $i$ is covered.
Therefore, the objective \eqref{eq:BP-obj} matches the reduced cost \eqref{eq:BP_redcost}.
So the IP finds a rule with minimum reduced cost.


Constraints \eqref{eq:BP-upper} can also be expressed as
\begin{equation}\label{eq:BP-misP-better}
    \delta_i + z_j \leq 1, \quad i \in \I, \; j \in S_i.
\end{equation}
This formulation may be favorable in terms of solver efficiency, as it provides a sharper LP relaxation.
These two formulations pose a tradeoff; while \eqref{eq:BP-misP-better} implies the \eqref{eq:BP-upper},
it is also a larger family of constraints which could significantly slow down the LP solver.
Furthermore, it is unclear how the pre-solvers employed modern IP solvers handle these formulations within this IP\@.
This makes it hard to predict which formulation is favorable.

%Here, $\I^{+}$ and $\I^{-}$ may simply be replaced by $\I$. \todo{explain why / why one would use se particular sets}.